% Modern editorial, markup and figure/layout contributions: CC-BY-SA-4.0.
% Historical text and illustrations: Leonhard Euler (1744), public domain.
% Component scope and upstream attribution: see RIGHTS.md.
% Caput IV. Printed pp.129–170; PDF132–173. Image-collated working transcription.
\subsection*{Caput IV.}
\textit{De Usu Methodi hactenus traditæ in resolutione varii generis quæstionum.}
\subsection*{Propositio I. Problema.}
\originalsection{1}\sourcepage{129}Invenire æquationem inter binas variabiles $x$ \& $y$, ita ut, pro dato ipsius $x$ valore, puta posito $x=a$, formula $\int Zdx$ obtineat maximum minimumve valorem, existente $Z$ functione ipsarum $x$, $y$, $p$, $q$, $r$, \&c. sive determinata sive indeterminata.
\subsection*{Solutio.}
Ex quacunque consideratione variabiles $x$ \& $y$ sint natæ, eæ \sourcepage{130}semper tanquam coordinatæ orthogonales cujuspiam curvæ considerari possunt: atque hanc ob rem quæstio proposita huc revocatur, ut determinetur curva abscissam habens $=x$ \& applicatam $=y$, in qua valor $\int Zdx$, si ad abscissam datæ magnitudinis, puta $x=a$, applicetur, fiat omnium maximus vel minimus. Quod si autem Problema hoc modo proponatur, tum ejus solutio in præcedentibus Capitibus satis superque est tradita. Quamobrem formulæ propositæ $\int Zdx$, secundum Methodos ante expositas, capi oportet valorem differentialem, qui datæ abscissæ $x=a$ conveniat, isque nihilo æqualis positus dabit æquationem inter $x$ \& $y$ desideratam, quæ pro data abscissa $x=a$, producet formulæ $\int Zdx$ maximum minimumve valorem. Q. E. I.
\subsection*{Corollarium I.}
\originalsection{2}Methodus ergo ante tradita multo latius patet, quam ad æquationes inter coordinatas curvarum inveniendas, ut quæpiam expressio $\int Zdx$ fiat maximum minimumve. Extenditur scilicet ad binas quascunque variabiles, sive eæ ad curvam aliquam pertineant quomodocunque, sive in sola analytica abstractione versentur.
\subsection*{Corollarium II.}
\originalsection{3}Inter binas autem variabiles propositas discrimen ingens intercedit, eo quod proposita formula $\int Zdx$ pro determinato quodam alterius variabilis valore maximum minimumve obtinere debeat. Isthanc ergo variabilem constanter litera $x$, alteram vero littera $y$ denotari convenit.
\subsection*{Corollarium III.}
\originalsection{4}Litteris igitur $x$ \& $y$ debito modo binis quantitatibus variabilibus inpositis, erit
\[p=\frac{dy}{dx},\quad q=\frac{dp}{dx},\quad r=\frac{dq}{dx},\quad s=\frac{dr}{dx}\quad\&c.\]
\sourcepage{131}His scilicet litteris differentialia cujuscunque gradus, quæ forte in maximi minimive formula insint, tolli poterunt, ita ut $Z$ futura sit functio litterarum $x$, $y$, $p$, $q$, $r$, \&c.
\subsection*{Corollarium IV.}
\originalsection{5}Cum ergo maximi minimive formula ad talem formam $\int Zdx$ fuerit reducta, in qua $Z$ sit functio ipsarum $x$, $y$, $p$, $q$, $r$, \&c. sive definita sive indefinita, tum ex superioribus præceptis formulæ $\int Zdx$ valor differentialis, respondens toti abscissæ propositæ $x=a$, debet investigari, qui nihilo æqualis positus præbebit æquationem inter $x$ \& $y$ quæsitam.
\subsection*{Corollarium V.}
\originalsection{6}Si $Z$ est functio definita ipsarum $x$, $y$, $p$, $q$, $r$, \&c. tum valor differentialis formulæ $\int Zdx$ non pendet a præscripto abscissæ valore $x=a$; \& hanc ob rem æquatio inter $x$ \& $y$ inventa pro qualibet abscissa præbebit maximum vel minimum formulæ $\int Zdx$.
\subsection*{Scholion I.}
\originalsection{7}Quia in hoc negotio valores differentiales, quos ante pro omni genere formularum sparsim eruimus, in promtu esse oportet, eos hic conjunctim in conspectum producemus, ut sit unde, quovis casu oblato, valores differentiales, quibus opus fuerit, conquiri ac depromi queant. Exhibebimus igitur formulæ $\int Zdx$ pro varia functionis $Z$ indole valorem differentialem, qui perpetuo determinatæ variabilis $x$ quantitati, puta $x=a$, respondeat.
\subsection*{I.}
\sourcepage{132}Maximi minimive formula $\int Zdx$.
\[dZ=Mdx+Ndy+Pdp+Qdq+Rdr+\&c.\]
Valor differentialis erit
\[nv.dx\left(N-\frac{dP}{dx}+\frac{ddQ}{dx^2}-\frac{d^3R}{dx^3}+\frac{d^4S}{dx^4}-\&c.\right)\]
qui valor differentialis pro omni variabilis $x$ magnitudine æque valet.
\subsection*{II.}
Maximi minimive formula $\int Zdx$.
\[dZ=Ld\Pi+Mdx+Ndy+Pdp+Qdq+\&c.\qquad\&\qquad\Pi=\int[Z]dx\]
existente
\[d[Z]=[M]dx+[N]dy+[P]dp+[Q]dq+[R]dr+\&c.\]
Jam posito post integrationem $x=a$, sit $\int Ldx=H$, ponaturque $H-\int Ldx=V$, Valor differentialis erit
\[nv.dx\left(\begin{aligned}N+[N]V&-\frac{d.(P+[P]V)}{dx}+\frac{dd.(Q+[Q]V)}{dx^2}\\&-\frac{d^3.(R+[R]V)}{dx^3}+\frac{d^4.(S+[S]V)}{dx^4}-\&c.\end{aligned}\right)\]
\subsection*{III.}
Maximi minimive formula $\int Zdx$.
\[\begin{aligned}dZ&=Ld\Pi+Mdx+Ndy+Pdp+Qdq+\&c.\\\Pi&=\int[Z]dx,\\d[Z]&=[L]d\pi+[M]dx+[N]dy+[P]dp+[Q]dq+\&c.\\\pi&=\int[z]dx,\\d[z]&=[m]dx+[n]dy+[p]dp+[q]dq+[r]dr+\&c.\end{aligned}\]
Sit iterum, posito post integrationem $x=a$ ut ante, $\int Ldx=H$
\sourcepage{133}ac ponatur $H-\int Ldx=V$. Jam integretur $\int[L]Vdx$, sitque integrale, eo casu quo $x=a$ ponitur, $=G$, ac ponatur
\[G-\int[L]Vdx=[V]=G-\int[L]dx\left(H-\int Ldx\right);\]
His positis, Valor differentialis erit
\[nv.dx\left(\begin{aligned}N+[N]V+[n][V]&-\frac{d.(P+[P]V+[p][V])}{dx}\\&+\frac{dd.(Q+[Q]V+[q][V])}{dx^2}\\&-\frac{d^3.(R+[R]V+[r][V])}{dx^3}\\&+\frac{d^4.(S+[S]V+[s][V])}{dx^4}-\&c.\end{aligned}\right)\]
unde simul lex progressionis patet, si adhuc plura integralia involvantur.
\subsection*{IV.}
Maximi minimive formula $\int Zdx$.
\[dZ=Ld\Pi+Mdx+Ndy+Pdp+Qdq+\&c.\qquad\&\qquad\Pi=\int Zdx\]
Abeat, posito $x=a$, hæc expressio $e^{\int Ldx}$ in $H$, denotante $e$ numerum cujus logarithmus est $=1$, sitque $He^{-\int Ldx}=V$: Valor differentialis erit
\[nv.dx\left(NV-\frac{d.PV}{dx}+\frac{dd.QV}{dx^2}-\frac{d^3RV}{dx^3}+\frac{d^4SV}{dx^4}-\&c.\right)\]
\subsection*{V.}
Maximi minimive formula $\int Zdx$.
\[\begin{aligned}dZ&=Ld\Pi+Mdx+Ndy+Pdp+Qdq+Rdr+\&c.\\\Pi&=\int[Z]dx,\\d[Z]&=[L]d\Pi+[M]dx+[N]dy+[P]dp+[Q]dq+[R]dr+\&c.\end{aligned}\]
Sit, si ponatur $x=a$, post integrationem
\[\int e^{\int[L]dx}Ldx=H\]
atque ponatur
\sourcepage{134}\[e^{-\int[L]dx}\left(H-\int e^{\int[L]dx}Ldx\right)=V\]
Valor differentialis erit
\[nv.dx\left(\begin{aligned}N+[N]V&-\frac{d.(P+[P]V)}{dx}+\frac{dd.(Q+[Q]V)}{dx^2}\\&-\frac{d^3.(R+[R]V)}{dx^3}+\frac{d^4.(S+[S]V)}{dx^4}-\&c.\end{aligned}\right).\]
In his igitur quinque casibus continentur omnes regulæ, quas in Capitibus præcedentibus invenimus. Iique tam late patent, ut omnes casus qui quidem occurrere queant, in iis vel actu contineantur, vel saltem per eos non difficulter resolvi possint. Iis igitur hic in compendium redactis, eorum usum monstrabimus, in resolvendis quæstionibus, in quibus $x$ \& $y$ non denotant coordinatas orthogonales.
\subsection*{Exemplum I.}
\originalsection{8}\marginnote{\footnotesize\hyperlink{fig:7}{Fig.~7.}}Ex dato centro $C$ ductis radiis $CA$, $CM$, invenire lineam $AM$, quæ inter omnes alias lineas intra angulum $ACM$ contentas sit brevissima.

Patet quidem hanc lineam quæsitam esse rectam: interim tamen hanc quæstionem secundum præcepta data resolvi conveniet, ut consensus Methodi cum veritate luculentius perspiciatur. Cum igitur longitudo lineæ $AM$ pro dato angulo $ACM$ debeat esse minima; ponamus angulum hunc $ACM$ esse $=x$; seu centro $C$, radio $CB=1$, describamus circulum, sitque arcus $BS=x$. Tum sit radius $CM$ altera variabilis $=y$, æquatione enim inter has variabiles $x$ \& $y$ inventa, innotescet natura lineæ quæsitæ $AM$. Jam autem ducto radio proximo $Cm$ erit $Ss=dx$, \& $mn=dy$, sumto $Cn=CM$: ob triangula vero similia $CSs$ \& $CMn$ erit $1:dx=CM[y]:Mn[ydx]$. Ex his itaque erit $Mm=\sqrt{dy^2+y^2dx^2}$; \& quia perpetuo ponimus $dy=pdx$, erit $Mm=dx\sqrt{yy+pp}$; unde lineæ $AM$ longitudo erit $=\int dx\sqrt{yy+pp}$, quæ debet esse minima pro dato ipsius $x$ valore, puta $x=a$. At quia hæc formula \sourcepage{135}ad casum primum pertinet, linea satisfaciens erit pro quovis valore ipsius $x$ minima. Cum igitur sit $Z=\sqrt{yy+pp}$, erit
\[dZ=\frac{ydy}{\sqrt{yy+pp}}+\frac{pdp}{\sqrt{yy+pp}},\]
\& in casu primo fiet
\[M=0,\quad N=\frac{y}{\sqrt{yy+pp}},\quad P=\frac{p}{\sqrt{yy+pp}},\quad Q=0,\quad R=0,\quad\&c.\]
ideoque $dZ=Ndy+Pdp$. Habebitur ergo iste valor differentialis $nv.dx(N-dP/dx)$, indeque pro solutione ista æquatio, $0=N-dP/dx$: quæ, multiplicata per $pdx=dy$, dat $Ndy=pdP$; quo in æquatione $dZ=Ndy+Pdp$ substituto, prodibit $dZ=Pdp+pdP$, \& integrando $Z+C=Pp$, seu
\[C+\sqrt{yy+pp}=\frac{pp}{\sqrt{yy+pp}}.\]
Quocirca erit $\dfrac{yy}{\sqrt{yy+pp}}=\mathit{Const.}=b$. At est
\[Mm[dx\sqrt{yy+pp}]:Mn[ydx]=MC[y]:\frac{yy}{\sqrt{yy+pp}};\]
quæ quarta proportionalis præbet perpendiculum $CP$, quod ex $C$ in tangentem lineæ quæsitæ $MP$ demittitur. Cum igitur hoc perpendiculum $CP$ sit constans, intelligitur lineam quæsitam esse rectam: \& quia, in æquatione inventa prima $Ndx=dP$, duæ insunt potentia constantes arbitrariæ, conditio hæc quæstioni est addenda, ut linea quæsita per data duo puncta transeat; tum igitur linea recta per illa duo puncta ducta quæsito satisfaciet.
\subsection*{Exemplum II.}
\originalsection{9}\marginnote{\footnotesize\hyperlink{fig:8}{Fig.~8.}}Super axe $AP$ construere lineam $BM$, ita comparatam, ut, abscissa area $ABMP$ datæ magnitudinis, arcus curvæ $BM$ illi areæ respondens sit omnium minimus.

Quia pro data area $ABMP$ minima longitudo arcus $BM$ requiritur, area $ABMP$ nobis designanda erit variabili $x$: altera variabili $y$ autem indicemus applicatam curvæ $PM$. Jam sit abscissa $AP=t$, erit $x=\int ydt$, ideoque $dt=dx/y$: atque \sourcepage{136}arcus $BM$ longitudo erit $=\int\sqrt{dy^2+dx^2/yy}$. Posito ergo $dy=pdx$, minimum esse debet hæc formula
\[\int dx\sqrt{\frac{1}{yy}+pp}=\int\frac{dx\sqrt{1+yypp}}{y}.\]
Erit itaque $Z=\dfrac{\sqrt{1+yypp}}{y}$, \&
\[dZ=-\frac{dy}{yy\sqrt{1+y^2p^2}}+\frac{yypdp}{y\sqrt{1+y^2p^2}};\]
unde $M=0$, $N=\dfrac{-1}{yy\sqrt{1+y^2p^2}}$; $P=\dfrac{yp}{\sqrt{1+y^2p^2}}$, $Q=0$ \&c. Pertinet ergo hæc quæstio ad casum primum, ac solutio præbebit lineam curvam, quæ pro area quacunque $APMB$ abscissa erit brevissima. Pervenietur autem, uti in præcedente Exemplo, ad æquationem hanc $Z=C+Pp$, atque curva quæsita per data duo puncta describi poterit. Erit itaque
\[\frac{\sqrt{1+yypp}}{y}=C+\frac{ypp}{\sqrt{1+yypp}},\]
seu $1=Cy\sqrt{1+yypp}$: vel $b=y\sqrt{1+yypp}$; hinc fit $bb=yy+y^4pp$, \&
\[p=\frac{\sqrt{bb-yy}}{yy}=\frac{dy}{dx}=\frac{dy}{ydt},\]
ob $dx=ydt$. Erit igitur $dt=\dfrac{ydy}{\sqrt{bb-yy}}$, \& $t=c\pm\sqrt{bb-yy}$. Quare linea quæsita erit Circulus, centro alicubi in axe $AP$, puta in $C$, assumto: isque inter omnes alias curvas per eadem duo quæcunque puncta ductas, pro data resecta area $ABMP$, habebit arcum $BM$ brevissimum.
\subsection*{Exemplum III.}
\originalsection{10}\marginnote{\footnotesize\hyperlink{fig:9}{Fig.~9.}}Eductis ex puncto fixo $C$ radiis $CA$, $CM$; intra eos describere curvam $AM$, quæ pro dato spatio $ACM$ habeat arcum $AM$ brevissimum.

Quia arcus $AM$ minimus esse debet, si spatium $ACM$ datæ magnitudinis abscindatur; ponatur area hæc $ACM=x$, atque radius $CM$ designetur altera variabili $y$. Jam ponatur
\sourcepage{137}arcus $BS$, radio $CB=1$ descriptus, $=t$; erit, ut ante vidimus, $Mn=ydt$, \& area $MCm=\tfrac12yydt=dx$, unde fit $dt=2dx/yy$. Quia porro est
\[Mm=\sqrt{dy^2+y^2dt^2}=\sqrt{dy^2+\frac{4dx^2}{yy}};\]
sit $dy=pdx$, minimumque esse debet $\int\dfrac{dx}{y}\sqrt{4+ppyy}$. Cum igitur sit $Z=\dfrac{\sqrt{4+y^2p^2}}{y}$, erit
\[M=0,\quad N=-\frac{4}{yy\sqrt{4+y^2p^2}},\quad\&\quad P=\frac{yp}{\sqrt{4+y^2p^2}},\quad Q=0,\quad\&c.\]
Hinc resultat ista æquatio $Z=C+Pp$; propterea quod sit $M=0$: ideoque
\[\frac{\sqrt{4+y^2p^2}}{y}=C+\frac{ypp}{\sqrt{4+y^2p^2}},\]
seu $4=Cy\sqrt{4+yypp}$ vel $2b=y\sqrt{4+yypp}$; hincque
\[p=\frac{2\sqrt{bb-yy}}{yy}=\frac{dy}{dx}=\frac{2dy}{yydt};\qquad\text{ac}\qquad dt=\frac{dy}{\sqrt{bb-yy}}:\]
itemque integrando
\[t=\mathrm{A\ sin.}\frac{y}{b}+\mathrm{A\ sin.}\frac{c}{b}=\mathrm{A\ sin.}\frac{y\sqrt{bb-cc}+c\sqrt{bb-yy}}{bb}.\]
In $AC$ ex $S$ demittatur perpendiculum $QS=\sin.\mathrm{A}t$, erit $QS=\dfrac{y\sqrt{bb-cc}+c\sqrt{bb-yy}}{bb}$. At ex æquatione $t+\mathit{Const.}=\mathrm{A\ sin.}\dfrac{y}{b}$ colligitur curva quæsita esse Circulus $AME$ per punctum fixum $C$ transiens. Describatur enim super diametro quacunque $CE$ in $C$ terminata Circulus $CAME$, arcus $AM$ interceptus inter radios $ACM$ pro data area $ACM$ erit minimus. Scilicet si alia quæcunque curva per duo quæcunque puncta in hoc Circulo sita describatur, binisque radiis ex $C$ ductis area æqualis areæ $ACM$ abscindatur, arcus illius curvæ respondens perpetuo major erit quam arcus $AM$. Quod ut appareat, ducatur ex $C$ ad $CE$ normalis $CD$, in eamque ex $S$ perpendiculum $SQ$ demittatur: erit triangulum $SCQ$ simile triangulo $CEM$, hincque $CE:CM[y]=CS[1]:SQ$ seu $SQ=\dfrac{y}{CE}=\sin.\mathrm{A}.DBS$, vel $DBS=\mathrm{A\ sin.}\dfrac{y}{CE}$. \sourcepage{138}Posita ergo diametro $CE=b$, \& quia est $DBS=BS+BD=t+\mathit{Const.}$ erit $t+\mathit{Const.}=\mathrm{A\ sin.}\dfrac{y}{b}$: quæ est ipsa illa proprietas, qua curvam quæsitam præditam esse oportere invenimus.
\subsection*{Exemplum IV.}
\originalsection{11}\marginnote{\footnotesize\hyperlink{fig:10}{Fig.~10.}}In superficie quacunque, sive convexa sive concava, ducere lineam, quæ sit intra suos terminos omnium brevissima.

Sumatur planum quodcunque ad quod superficies referatur, $APQ$, in eoque capiatur recta $AP$ pro axe. Jam ex lineæ quæsitæ singulis punctis concipiantur perpendicula in hoc planum demitti, quibus describatur linea $AQ$, quæ erit projectio lineæ brevissimæ in hoc planum; qua cognita, simul ipsa linea brevissima in superficie proposita innotescet. Vocetur $AP=x$, $PQ=y$; atque cum natura superficiei detur, ex datis $AP=x$ \& $PM=y$ definiri poterit longitudo perpendicularis $QM$ in planum $APQ$, donec superficiem in $M$ secet. Quod si ergo ponatur $QM=z$, longitudo hujus lineæ $z$ dabitur per $x$ \& $y$, ita ut $z$ sit functio definita ipsarum $x$ \& $y$. Cum igitur sit $z$ functio ipsarum $x$ \& $y$, quæ ex æquatione locali ad superficiem datur, ponamus esse $dz=Tdx+Vdy$; eruntque $T$ \& $V$ ejusmodi functiones ipsarum $x$ \& $y$, ut $Tdx+Vdy$ sit formula differentialis definita: posito nempe $dT=Edx+Fdy$, erit $dV=Fdx+Gdy$, existente littera $F$ utrique differentiali communi. Nunc elementum lineæ in superficie ductæ est
\[=\sqrt{dx^2+dy^2+dz^2}=\sqrt{dx^2+dy^2+(Tdx+Vdy)^2}.\]
Posito ergo $dy=pdx$, minimum esse debet hæc formula $\int dx\sqrt{1+p^2+T^2+2TVp+V^2p^2}$; ita ut sit $Z=\sqrt{1+p^2+T^2+2TVp+V^2p^2}$, unde fit
\sourcepage{139}\[
dZ=\frac{\begin{aligned}&TEdx+TFdy+pdp\\&+VEpdx+VF pdy+TVdp\\&+TFpdx+TGpdy+V^2pdp\\&+VF p^2dx+VGppdy\end{aligned}}{\sqrt{1+pp+T^2+2TVp+V^2p^2}}.
\]
Quæ formula cum ad casum primum pertineat, proveniet ista æquatio inter $x$ \& $y$;
\[\frac{TFdx+VF pdx+TGpdx+VGppdx}{\sqrt{1+pp+T^2+2TVp+V^2p^2}}=d.\frac{p+TV+V^2p}{\sqrt{1+pp+T^2+2TVp+V^2p^2}}.\]
Est vero $Fdx+Gpdx=Fdx+Gdy=dV$, unde erit
\[\begin{aligned}\frac{TdV+VpdV}{\sqrt{1+pp(T+Vp)^2}}&=d.\frac{p+TV+V^2p}{\sqrt{1+pp+(T+Vp)^2}}\\&=\frac{\begin{aligned}&dp(1+T^2+V^2)+dT(V-Tp)\\&+dV(T+T^3+3T^2Vp+3TV^2p^2+V^3p^3+2Vp+Vp^3)\end{aligned}}{(1+pp+(T+Vp)^2)^{3:2}}.\end{aligned}\]
Æquatione autem ordinata, resultabit hæc
\[dp(1+T^2+V^2)+dT(V-Tp)+dV(Vp-Tpp)=0,\]
seu $dp=\dfrac{(Tp-V)(dT+pdV)}{1+T^2+V^2}$. Cum vero sit $p=dy/dx$, erit $dp=ddy/dx$; hincque fiet
\[dxddy=\frac{(Tdy-Vdx)(dxdT+dydV)}{1+T^2+V^2}\]
quæ est æquatio differentio-differentialis pro projectione $AQ$ lineæ brevissimæ in superficie quæsita; ideoque indicat, eam per duo quæque puncta duci posse. Æquatio hæc inventa in varias formas transmutari potest, quæ sæpius majore commodo usurpari poterunt. Ac primo quidem expedit eliminari differentialia $dT$ \& $dV$: cum enim sit $dz=Tdx+Vdy$, erit $ddz=dxdT+dydV+Vddy$; ideoque $dxdT+dydV=ddz-Vddy$, \sourcepage{140}quo valore substituto prodibit ista æquatio
\[dxddy+T^2dxddy+V^2dxddy=Tdyddz-Vdxddz-TVdyddy+V^2dxddy,\]
seu $dxddy+Tdzddy=Tdyddz-Vdxddz$; hincque $ddy:ddz=Tdy-Vdx:dx+Tdz$. Multiplicetur æquatio inventa per $dz$, ac in primo termino scribatur $Tdx+Vdy$ loco $dz$, erit
\[Tdx^2ddy+Vdxdyddy+Tdz^2ddy=Tdydzddz-Vdxdzddz.\]
Addatur utrinque $Tdy^2ddy-Vdxdyddy$, erit
\[Tddy(dx^2+dy^2+dz^2)=(dzddz+dyddy)(Tdy-Vdx)\]
seu
\[\frac{dyddy+dzddz}{dx^2+dy^2+dz^2}=\frac{Tddy}{Tdy-Vdx}=\frac{Tddz}{dx+Tdz}.\]
Vel multiplicetur æquatio per $dx$, ac loco $Tdx$ scribatur $dz-Vdy$, obtinebitur
\[dx^2ddy+dz^2ddy-Vdydzddy=dydzddz-Vdy^2ddz-Vdx^2ddz.\]
Addatur utrinque $dy^2ddy-Vdz^2ddz$, erit
\[\begin{aligned}ddy(dx^2+dy^2+dz^2)-Vdz(dyddy+dzddz)&=dy(dyddy+dzddz)\\&\quad-Vddz(dx^2+dy^2+dz^2);\end{aligned}\]
ideoque
\[\frac{dyddy+dzddz}{dx^2+dy^2+dz^2}=\frac{ddy+Vddz}{dy+Vdz}:\]
quæ æquationes omnes in sequenti expressione continentur
\[\frac{dyddy+dzddz}{dx^2+dy^2+dz^2}=\frac{Tddy}{Tdy-Vdx}=\frac{Tddz}{dx+Tdz}=\frac{ddy+Vddz}{dy+Vdz}.\]
Hic notandum est, quia quantitatum $T$ \& $V$ differentialia nusquam occurrunt, perinde esse, sive in $T$ \& $V$ contineatur $z$, sive minus. Quovis igitur casu oblato, conveniet eam æquationem assumere, quæ facillime integrationem admittat. Veluti si superficies proposita sit solidi rotundi conversione cujuscunque figuræ circa axem $AP$ nati, erit $yy+zz=$ quadrato functionis ipsius $x$, quæ sit $=X$, estque applicata illius curvæ genitricis abscissæ $x$ respondens. Erit itaque $zdz=XdX-ydy$, \& $dz=\dfrac{XdX}{z}-\dfrac{ydy}{z}$ unde fiet $T=\dfrac{XdX}{zdx}$, \& $V=-y/z$. Sumatur jam, commodi ergo, æquatio in qua $T$ non
\sourcepage{141}occurrit, hæc
\[\frac{dyddy+dzddz}{dx^2+dy^2+dz^2}=\frac{ddy+Vddz}{dy+Vdz},\]
quæ, ob $V=-y/z$, transit in hanc
\[\frac{dyddy+dzddz}{dx^2+dy^2+dz^2}=\frac{zddy-yddz}{zdy-ydz},\]
cujus integrale est $l\sqrt{dx^2+dy^2+dz^2}=l\dfrac{zdy-ydz}{b}$, seu $zdy-ydz=b\sqrt{dx^2+dy^2+dz^2}$. Quoniam nunc est $z=\sqrt{X^2-y^2}$, ponatur $dX=vdx$, erit $dz=\dfrac{Xvdx-ydy}{\sqrt{X^2-y^2}}$, \&
\[\begin{aligned}zdy-ydz&=\frac{X^2dy-Xyvdx}{\sqrt{X^2-y^2}},\\\sqrt{dx^2+dy^2+dz^2}&=\frac{\sqrt{X^2dx^2-y^2dx^2+X^2dy^2+X^2v^2dx^2-2Xyvdxdy}}{\sqrt{X^2-y^2}}.\end{aligned}\]
Ergo
\[\begin{aligned}X^4dy^2-2X^3yvdxdy+X^2y^2v^2dx^2&=bbX^2dx^2-bby^2dx^2+bbX^2dy^2\\&\quad+b^2X^2v^2dx^2-2b^2Xyvdxdy,\end{aligned}\]
seu
\[dy^2=\frac{2(b^2-X^2)Xyvdxdy+X^2y^2v^2dx^2-b^2X^2dx^2+b^2y^2dx^2-b^2X^2v^2dx^2}{X^2(bb-XX)},\]
quæ, extracta radice, præbet
\[dy=\frac{yvdx}{X}\pm\frac{bdx\sqrt{(1+vv)(yy-XX)}}{X\sqrt{bb-XX}}.\]
Quod si ponatur $y=Xt$, ut sit $dy=Xdt+tvdx$, fiet
\[\frac{dt}{\sqrt{tt-1}}=\frac{bdx\sqrt{1+vv}}{X\sqrt{bb-XX}};\]
in qua æquatione, quia $X$ \& $v$ sunt functiones ipsius $x$, variabiles $t$ \& $x$ a se invicem sunt separatæ.
\subsection*{Exemplum V.}
\originalsection{12}\marginnote{\footnotesize\hyperlink{fig:11}{Fig.~11.}}Super axe $APN$ construere curvam $AM$ ejusmodi, ut, abscissa per normalem $MN$ area $ANM$ datæ magnitudinis, arcus $AM$ sit minimus.

Quia, pro definita areæ $AMN$ magnitudine, arcus $AM$ minimus esse debet, ponatur area $AMN=ax$, positoque $x=a$, quo casu area $AMN$ fit $=aa$, fiat arcus $AM$ minimus. Ponatur porro applicata orthogonalis $MP=y$, abscissa $AP=t$, \& subnormalis $PN=u$; erit $ax=\int ydt+\tfrac12uy$, \& $u=\dfrac{ydy}{dt}$: elementum vero arcus $AM$ erit $=\dfrac{dy\sqrt{yy+uu}}{u}$. \sourcepage{142}Porro cum sit $adx=ydt+\tfrac12(udy+ydu)$ \& $dt=ydy/u$, erit $audx=yydy+\tfrac12uudy+\tfrac12yudu$, \& $du=\dfrac{2adx}{y}-\dfrac{2ydy}{u}-\dfrac{udy}{y}$. Jam ponatur $dy=pdx$, minimum esse debebit $\int\dfrac{pdx\sqrt{yy+uu}}{u}$, atque $u$ est quantitas cujus valor ex hac æquatione $du=dx\left(\dfrac{2a}{y}-\dfrac{2yp}{u}-\dfrac{up}{y}\right)$ definiri debet. Pertinet itaque hæc quæstio ad Casum quintum; cum quo si comparatio instituatur, fit $u=\Pi$ \& $Z=\dfrac{p\sqrt{yy+\Pi^2}}{\Pi}$, unde
\[L=\frac{-pyy}{\Pi^2\sqrt{yy+\Pi^2}};\quad M=0,\quad N=\frac{yp}{\Pi\sqrt{yy+\Pi^2}},\quad\&\quad P=\frac{\sqrt{yy+\Pi\Pi}}{\Pi}.\]
Deinde cum sit $\Pi=\int dx\left(\dfrac{2a}{y}-\dfrac{2yp}{\Pi}-\dfrac{\Pi p}{y}\right)$; fit $[Z]=\dfrac{2a}{y}-\dfrac{2yp}{\Pi}-\dfrac{\Pi p}{y}$, \& differentiando erit
\[[L]=\frac{2yp}{\Pi^2}-\frac{p}{y};\quad[M]=0,\quad[N]=\frac{-2a}{yy}-\frac{2p}{\Pi}+\frac{\Pi p}{yy}\quad\&\quad[P]=\frac{-2y}{\Pi}-\frac{\Pi}{y}.\]
Jam erit $\int[L]dx=\int\dfrac{2ydy}{\Pi^2}-ly$, \& $e^{\int[L]dx}=\dfrac{e^{\int 2ydy:\Pi\Pi}}{y}$; at est $Ldx=\dfrac{-yydy}{\Pi^2\sqrt{yy+\Pi\Pi}}$; unde fiet
\[\int e^{\int[L]dx}Ldx=-\int\frac{e^{\int 2ydy:\Pi\Pi}ydy}{\Pi^2\sqrt{yy+\Pi^2}},\]
cujus valor posito $x=a$, fiat $=H$, sitque
\[V=e^{-\int 2ydy:\Pi\Pi}y\left(H+\int\frac{e^{\int 2ydy:\Pi\Pi}ydy}{\Pi^2\sqrt{yy+\Pi^2}}\right).\]
His præparatis, erit æquatio satisfaciens $(N+[N]V)dx=d.(P+[P]V)$, sive substitutionibus factis,
\[\frac{ydy}{\Pi\sqrt{yy+\Pi\Pi}}-\frac{2aVdx}{yy}-\frac{2Vdy}{\Pi}-\frac{\Pi Vdy}{yy}=d.\left(\frac{\sqrt{yy+\Pi\Pi}}{\Pi}-\frac{2Vy}{\Pi}-\frac{\Pi V}{y}\right).\]
At est $2adx=yd\Pi+\dfrac{2yydy}{\Pi}+\Pi dy$; \sourcepage{143}unde erit
\[\begin{aligned}\frac{ydy}{\Pi\sqrt{yy+\Pi^2}}-\frac{Vd\Pi}{y}-\frac{4Vdy}{\Pi}
&=d.\left(\frac{\sqrt{yy+\Pi^2}}{\Pi}-\frac{2Vy}{\Pi}-\frac{\Pi V}{y}\right)\\
&=\frac{ydy}{\Pi\sqrt{yy+\Pi^2}}-\frac{yyd\Pi}{\Pi^2\sqrt{yy+\Pi^2}}\\&\quad-\frac{2Vdy}{\Pi}-\frac{2ydV}{\Pi}+\frac{2Vyd\Pi}{\Pi^2}\\&\quad-\frac{\Pi dV}{y}-\frac{Vd\Pi}{y}+\frac{\Pi Vdy}{yy};\end{aligned}\]
hincque
\[\frac{yyd\Pi}{\Pi^2\sqrt{y^2+\Pi^2}}-\frac{2Vdy}{\Pi}+\frac{2ydV}{\Pi}-\frac{2Vyd\Pi}{\Pi^2}+\frac{\Pi dV}{y}-\frac{\Pi Vdy}{yy}=0.\]
Verum, est generaliter $dV=-Ldx-V[L]dx$; unde erit
\[dV=\frac{yydy}{\Pi^2\sqrt{yy+\Pi\Pi}}-\frac{2Vydy}{\Pi^2}+\frac{Vdy}{y};\]
hincque
\[\frac{yyd\Pi}{\Pi^2\sqrt{yy+\Pi\Pi}}=\frac{dVd\Pi}{dy}+\frac{2Vyd\Pi}{\Pi^2}-\frac{Vd\Pi}{y};\]
quo substituto oritur
\[\frac{dVd\Pi}{dy}-\frac{2Vdy}{\Pi}+\frac{2ydV}{\Pi}+\frac{\Pi dV}{y}-\frac{Vd\Pi}{y}-\frac{\Pi Vdy}{yy}=0;\]
hoc est
\[\begin{aligned}dV\left(\frac{d\Pi}{dy}+\frac{2y}{\Pi}+\frac{\Pi}{y}\right)&=V\left(\frac{d\Pi}{y}+\frac{2dy}{\Pi}+\frac{\Pi dy}{yy}\right)\\&=\frac{ydV}{dy}\left(\frac{d\Pi}{y}+\frac{2dy}{\Pi}+\frac{\Pi dy}{yy}\right);\end{aligned}\]
quæ æquatio, cum sit divisibilis per $\dfrac{d\Pi}{y}+\dfrac{2dy}{\Pi}+\dfrac{\Pi dy}{yy}$, duplicem dat solutionem. Quarum prima erit $dV/V=dy/y$, quæ præbet $V=cy$: quoniam vero $V$ evanescere debet in casu minimi, eodem casu erit $y=0$; scilicet posito $x=a$ fiet $y=0$. Cum nunc sit $V=cy$, facta substitutione in æquatione
\[dV=\frac{yydy}{\Pi^2\sqrt{y^2+\Pi^2}}-\frac{2Vydy}{\Pi^2}+\frac{Vdy}{y},\]
erit $\dfrac{yydy}{\Pi^2\sqrt{y^2+\Pi^2}}=\dfrac{2cyydy}{\Pi^2}$; hincque vel $y=0$, vel $dy=0$, quo casu prodit linea recta axi parallela; vel $\Pi=\infty$, quo casu prodit linea recta ad axem normalis: vel etiam $\sqrt{yy+\Pi\Pi}=MN=\mathit{Const.}$ quæ æquatio dat Circulum; atque integer semicirculus, ob $y=0$ in casu minimi, quæsito satisfacit. Secunda solutio prodit ex divisore $\dfrac{d\Pi}{y}+\dfrac{2dy}{\Pi}+\dfrac{\Pi dy}{yy}=0$, seu $\Pi d\Pi+\dfrac{\Pi\Pi dy}{y}+2ydy=0$, \sourcepage{144}quæ multiplicata per $yy$, fit $yy\Pi d\Pi+\Pi\Pi ydy+2y^3dy=0$, cujus integrale est $\Pi^2y^2+y^4=C$, hincque $\Pi=\dfrac{\sqrt{b^4-y^4}}{y}$; quæ æquatio, quia non pendet ab $V$, pro quocunque valore ipsius $x$ satisfaciet. Erit autem introducta abscissa $AP=t$, ob $u=\Pi=\dfrac{ydy}{dt}$, ista æquatio $\dfrac{ydy}{dt}=\dfrac{\sqrt{b^4-y^4}}{y}$, unde $dt=\dfrac{yydy}{\sqrt{b^4-y^4}}$, ex qua æquatione intelligitur Elasticam rectangulam quæsito satisfacere; ita ut pro area $ANM$ inter normales $AN$ \& $MN$ arcus curvæ $AM$ sit brevissimus. Hæc autem curva per data duo puncta, siquidem axis $AP$ sit positione datus, describi potest.
\subsection*{Scholion II.}
\originalsection{13}Ex his Exemplis eximius usus, quem habet nostra Methodus in Problematis etiam diversi generis resolvendis, abunde patet; inprimis autem ultimum Exemplum nonnullas notatu maxime dignas suppeditat circumstantias, ex quibus natura solutionis illustrari poterit. Quoniam enim duplex æquatio ob factores duos nata est, duplex quoque solutio prodiit; quarum prior lineam satisfacientem absolute determinat, ita ut ea per data duo puncta duci nequeat: dat enim vel lineam rectam, vel semicirculum. Linea recta duplici modo quæstionem solvit, dum est vel normalis ad axem $AP$, vel eidem parallela; \& quemadmodum utraque satisfaciat manifestum est: nam in ea, quæ est normalis ad axem, portio quæ cum axe \& normali datum spatium comprehendit perpetuo est infinite parva, ideoque revera minima: altera recta axi parallela aliquanto latius patet, cum ea per datum punctum duci possit; \& quia ipsæ applicatæ ad eam sunt normales, ac spatium abscissum sit ut ipsa abscissa, ejus respectu linea illa recta utique erit brevissima. Semicirculus deinde, qui ex prima solutione prodiit, ita absolute satisfacit, ut, proposita spatii abscindendi quantitate, ipse semicirculus determinetur,
\sourcepage{145}ejus enim area esse debet $=aa$. Secunda autem solutio, quæ curvam Elasticam rectangulam præbuit, latius patet: nam per data duo quæcunque puncta ejusmodi curva traduci potest, eaque, inter omnes alias curvas per eadem puncta transeuntes, hac gaudebit prærogativa, ut si, in omnibus curvis, per normales, areæ æquales abscindantur, arcus Elasticæ futurus sit omnium minimus. His igitur expositis pergamus ad usum Methodi traditæ ostendendum, in iis maximi minimive investigationibus, in quibus maximi minimive formula non est talis expressio integralis simplex $\int Zdx$, qualem formam hactenus perpetuo tractavimus; verum est composita ex duabus pluribusve hujusmodi formulis quomodocunque. Ac primo quidem, si maximum minimumve esse debeat aggregatum duarum pluriumve formularum integralium, puta $\int Zdx+\int Ydx-\int Xdx$, operatio nulla difficultate laborat: quia enim formula maximi minimive est $\int dx(Z+Y-X)$, hæc tanquam simplex formula integralis tractari, ejusque valor differentialis assignari poterit. Operatio autem eo redibit, ut pro singulis formulis $\int Zdx$, $\int Ydx$ \& $\int Xdx$, earum valores differentiales quærantur; earumque loco in formula $\int Zdx+\int Ydx-\int Xdx$ substituantur; \& quod oritur nihilo æquale ponatur: sicque habebitur æquatio quæsito satisfaciens.
\subsection*{Propositio II. Problema.}
\originalsection{14}Invenire æquationem inter $x$ \& $y$, ut, posito $x=a$, fiat hæc expressio $\int Zdx\times\int Ydx$, quæ est productum ex duabus formulis integralibus $\int Zdx$ \& $\int Ydx$, maximum vel minimum.
\subsection*{Solutio.}
Ponamus istam æquationem inter $x$ \& $y$ jam esse inventam, foreque ex ea, posito $x=a$, valorem Formulæ $\int Zdx=A$, \& $\int Ydx=B$; erunt hæ quantitates $A$ \& $B$ constantes; atque earum productum $AB$ maximum vel minimum. Jam ponatur apud valorem indefinitum $x$ variabilem $y$ augeri particula $nv$, \sourcepage{146}ex ea utraque quantitas $A$ \& $B$ incrementum accipiet, unaquæque scilicet augebitur valore differentiali ex præcedentibus definiendo. Sit igitur $dA$ valor differentialis ipsius $A$, qui respondet formulæ integrali $\int Zdx$, posito $x=a$, similique modo sit $dB$ valor differentialis ipsius $B$ oriundus ex formula $\int Ydx$, posito $x=a$. Cum ergo, ex adjecta particula $nv$ variabili $y$, abeat $A$ in $A+dA$, \& $B$ in $B+dB$, productum $AB$ transmutabitur in $AB+AdB+BdA+dAdB$; quare cum $AB$ esse debeat maximum vel minimum, oportebit esse $AB=AB+AdB+BdA+dAdB$. Ideoque $0=AdB+BdA$, ob evanescentem terminum $dAdB$ præ reliquis. Ex his itaque oritur sequens Problematis solutio; Quæratur formulæ $\int Zdx$ valor differentialis qui sit $dA$, sitque $A$ valor formulæ $\int Zdx$, quem obtinet posito $x=a$. Deinde quæratur formulæ $\int Ydx$ valor differentialis, qui sit $dB$, ac $B$ denotet valorem formulæ $\int Ydx$, quem recipit posito $x=a$: quibus factis habebitur ista æquatio $0=AdB+BdA$, in qua relatio satisfaciens inter $x$ \& $y$ continebitur. Q. E. I.
\subsection*{Corollarium I.}
\originalsection{15}Quanquam in æquatione $0=AdB+BdA$ insunt quantitates constantes $A$ \& $B$, tamen eæ non sunt arbitrariæ, sed utraque per ipsam hanc æquationem definietur. Scilicet si ex hac æquatione eliciantur valores $\int Zdx$ \& $\int Ydx$, ponaturque $x=a$, prodire debent illæ quantitates $A$ \& $B$; unde hæ determinabuntur per $a$, \& per reliquas constantes arbitrarias quæ per integrationem ingredientur.
\subsection*{Corollarium II.}
\originalsection{16}Si $Z$ \& $Y$ fuerint functiones determinatæ quantitatum $x$, $y$, $p$, $q$, $r$, \&c. tum valores differentiales $dA$ \& $dB$ non pendebunt ab $a$; interim tamen quantitas $a$ ingreditur in æquationem $0=AdB+BdA$: ex quo curva inventa, tantum pro definito abscissæ $x$ valore $x=a$, quæsito satisfaciet.
\subsection*{Corollarium III.}
\originalsection{17}\sourcepage{147}Ex æquatione autem $0=AdB+BdA$ particula $nv$ omnino egreditur: nam quia uterque valor differentialis $dA$ \& $dB$ per $nv$ multiplicatus prodiit, iterum $nv$ per divisionem exterminabitur: hocque modo æquatio inter $x$ \& $y$ atque constantes nascetur, quæ Problemati satisfiet.
\subsection*{Scholion I.}
\originalsection{18}Neminem hic forma æquationis $0=AdB+BdA$ inventæ offendat, eo quod speciem formulæ differentialis definitæ præ se ferat, neque hinc etiam quisquam concludat æquationis $0=AdB+BdA$ integralem sumi posse hanc, $\mathit{Const.}=AB$. Jam enim significationes explicavimus, quas tribuimus cum litteris $A$ \& $B$, tum etiam formis differentialibus $dA$ \& $dB$: ex quo intelligere licet, vulgarem notandi modum hic non locum habere. Ideo autem hunc notandi modum, etsi a consueto dissentientem, hic adhibere visum est, ut nexus æquationis $0=AdB+BdA$ cum formula maximi minimive $\int Zdx.\int Ydx$ melius perspiciatur. Cum enim maximum minimumve respondere debeat valori $x=a$; ponamus hoc casu abire $\int Zdx$ in $A$ \& $\int Ydx$ in $B$; quo facto, maximum minimumve erit $AB$. Hinc autem sponte nascitur æquatio inventa $0=AdB+BdA$, siquidem $AB$, litteris $A$ \& $B$ tanquam variabilibus spectatis, differentietur. Quod cum fuerit factum, in memoriam revocari oportet, pro differentialibus $dA$ \& $dB$ accipiendos esse valores differentiales eos, qui conveniunt formulis integralibus $\int Zdx$ \& $\int Ydx$, ex quibus ipsæ quantitates $A$ \& $B$ constantes prodiere. Hunc nexum ideo annotasse juvabit, quod infra eundem ad quemcunque compositionis modum, quo formula maximi minimive ex formulis integralibus composita fuerit, æque patere; similique modo ex ipsa maximi minimive expressione per differentiationem æquationem quæsitam obtineri ostendemus.
\subsection*{Exemplum I.}
\originalsection{19}\sourcepage{148}Invenire æquationem inter $x$ \& $y$, ut, posito $x=a$, fiat ista expressio $\int ydx\times\int xdy$ maximum.

Fiat $\int ydx=A$, \& $\int xdy=B$, posito $x=a$, \& quærantur formularum $\int ydx$ \& $\int xdy$, seu $\int xpdx$, valores differentiales: ac formulæ $\int ydx$ valor differentialis est $nv.dx.1$, formulæ autem $\int xdy$, seu $\int xpdx$, est $nv.dx\left(-\dfrac{1}{dx}d.x\right)=-nv.dx$. Erit ergo $dA=nv.dx$, \& $dB=-nv.dx$: unde æquatio $0=AdB+BdA$ abibit in hanc $0=-A.nv.dx+B.nv.dx$, seu $A=B$. Quæsito ergo omnes æquationes inter $x$ \& $y$ æque satisfaciunt, dummodo, casu $x=a$, fuerit $\int ydx=\int xdy$; hoc est area curvæ $=\tfrac12xy$.
\subsection*{Exemplum II.}
\originalsection{20}Invenire æquationem inter $x$ \& $y$, ut, casu $x=a$, fiat minimum hæc expressio $\int ydx\times\int dx\sqrt{1+pp}$.

Casu $x=a$, fiat $\int ydx=A$, \& $\int dx\sqrt{1+pp}=B$. Porro sumendis valoribus differentialibus erit $dA=nv.dx.1$, \&
\[dB=nv.dx\left(-\frac{1}{dx}d.\frac{p}{\sqrt{1+pp}}\right)=-nv.d.\frac{p}{\sqrt{1+pp}}.\]
Hinc prodit sequens æquatio
\[0=-A.nv.d.\frac{p}{\sqrt{1+pp}}+B.nv.dx,\qquad\text{seu}\qquad Bdx=Ad.\frac{p}{\sqrt{1+pp}}.\]
Quæ integrata dat $x+b=\dfrac{Ap}{B\sqrt{1+pp}}$, ubi $A/B$ denotat rationem, quam tenet $\int ydx$ ad $\int dx\sqrt{1+pp}$ tum cum sit $x=a$. Sit brevitatis gratia $A/B=c$, erit $(x+b)\sqrt{1+pp}=cp$, \& $p=\dfrac{x+b}{\sqrt{cc-(x+b)^2}}=dy/dx$. Integrata ergo hac æquatione,
\sourcepage{149}resultabit $y=f\pm\sqrt{cc-(x+b)^2}$, ita ut sit $(y-f)^2+(x+b)^2=c^2$, unde patet curvam satisfacientem esse Circulum, radio $c$ descriptum, axe ubicunque accepto. Hujusmodi vero Circuli non quivis arcus satisfaciet, verum is tantum qui per $c$ radium Circuli multiplicatus producit aream; est enim $A=Bc$. Ergo vel radius Circuli $c$ pro lubitu accipi potest, ex eoque definietur illa abscissæ $x$ magnitudo determinata $a$; vel si $a$ detur, ut ponimus, inde vicissim radius $c$ determinabitur. Perspicuum autem est arcum Circuli, qui satisfacit, convexitate sua axem respicere debere; hoc enim casu area fit minor, ideoque productum ex area in arcum minimum.
\subsection*{Exemplum III.}
\originalsection{21}Invenire curvam, in qua, pro data abscissa $x=a$, minimum fiat hæc expressio $\int yxdx\times\int xdx\sqrt{1+pp}$.

Posito $x=a$, fiat $\int yxdx=A$, \& $\int xdx\sqrt{1+pp}=B$. Erit autem $dA=nv.dx.x$ \& $dB=-nv.dx.\dfrac{1}{dx}d.\dfrac{xp}{\sqrt{1+pp}}$; unde obtinebitur ista æquatio $Bxdx=Ad.\dfrac{px}{\sqrt{1+pp}}$, quæ integrata dat
\[xx\pm bb=\frac{2Apx}{B\sqrt{1+pp}}=\frac{2cpx}{\sqrt{1+pp}},\]
posito $A/B=c$. Hinc $p=\dfrac{xx\pm bb}{\sqrt{4ccxx-(xx\pm bb)^2}}=dy/dx$, ideoque pro curva habebitur hæc æquatio,
\[y=\int\frac{(xx\pm bb)dx}{\sqrt{4ccxx-(xx\pm bb)^2}}.\]
De qua notandum est, si fiat $b=0$, tum prodire æquationem pro Circulo $y=\int\dfrac{xdx}{\sqrt{4cc-xx}}$ cujus radius sit $2c$.
\subsection*{Scholion II.}
\originalsection{22}\sourcepage{150}Eadem hæc Exempla omnia quoque resolvi possunt per Methodum supra jam traditam; quare cum utraque via eadem solutio obtineatur, juvabit solutionem per alteram viam uno Exemplo exhiberi. Sumamus igitur tertium Exemplum, in quo maximi minimive formula $\int yxdx\times\int xdx\sqrt{1+pp}$, differentiando iterumque integrando per partes, reducitur ad hanc formam
\[\int yxdx\int xdx\sqrt{1+pp}+\int xdx\sqrt{1+pp}\int yxdx;\]
cujus utrumque membrum in Casu secundo supra §. 7. exposito continetur. Quæratur itaque utriusque valor differentialis, eorum enim summa, posita $=0$, dabit æquationem pro curva quæsita. Formula autem $\int yxdx\int xdx\sqrt{1+pp}$ cum Casu secundo collata, dabit $\Pi=\int xdx\sqrt{1+pp}$ \& $Z=yx\Pi$; unde fit $L=yx$; $M=y\Pi$, $N=x\Pi$, $P=0$, \&c. Deinde erit $[Z]=x\sqrt{1+pp}$; indeque $[M]=\sqrt{1+pp}$, $[N]=0$, \& $[P]=\dfrac{xp}{\sqrt{1+pp}}$. Porro est $\int Ldx=\int yxdx$, cujus valor, posito $x=a$, quem generaliter posuimus $H$, hic in solutione Exempli est $A$; ita ut sit $V=A-\int yxdx$. Quare hujus formulæ valor differentialis erit
\[\begin{aligned}&=nv.dx\left(x\Pi-\frac{1}{dx}d.\frac{xp(A-\int yxdx)}{\sqrt{1+pp}}\right)\\&=nv.dx\left(x\int xdx\sqrt{1+pp}-\frac{A}{dx}d.\frac{xp}{\sqrt{1+pp}}+\frac{1}{dx}d.\frac{xp\int yxdx}{\sqrt{1+pp}}\right).\end{aligned}\]
Altera formula $\int xdx\sqrt{1+pp}\int yxdx$, cum Casu secundo §. 7. collata, dat $\Pi=\int yxdx$ \& $Z=x\Pi\sqrt{1+pp}$, unde erit $L=x\sqrt{1+pp}$, $M=\Pi\sqrt{1+pp}$, $N=0$, \& $P=\dfrac{x\Pi p}{\sqrt{1+pp}}$; hincque $\int Ldx=\int xdx\sqrt{1+pp}$: quare cum $H$ sit valor ipsius $\int Ldx$, posito $x=a$, erit $H=B$, \& $V=B-\int xdx\sqrt{1+pp}$. Porro est $[Z]=yx$, hincque $[M]=y$, $[N]=x$, \& $[P]=0$. Ex his prodit valor differentialis
\[=nv.dx\left(Bx-x\int xdx\sqrt{1+pp}-\frac{1}{dx}\times d.\frac{xp\int yxdx}{\sqrt{1+pp}}\right).\]
\sourcepage{151}His igitur valoribus differentialibus ambobus additis, emerget hujus expressionis compositæ
\[\int yxdx\int xdx\sqrt{1+pp}+\int xdx\sqrt{1+pp}\int yxdx,\]
seu hujus $\int yxdx\times\int xdx\sqrt{1+pp}$, quæ in Exemplo erat proposita, valor differentialis
\[=nv.dx\left(Bx-\frac{A}{dx}d.\frac{xp}{\sqrt{1+pp}}\right),\]
ex quo pro curva æquatio erit hæc $Bxdx=Ad.\dfrac{xp}{\sqrt{1+pp}}$, quam eandem in solutione Exempli invenimus. Similis autem consensus in genere deprehendetur, si quis expressionem $\int Zdx\times\int Ydx$ eodem modo tractare voluerit.
\subsection*{Propositio III. Problema.}
\originalsection{23}Invenire æquationem inter $x$ \& $y$ ejus conditionis, ut, posito $x=a$, ista fractio $\dfrac{\int Zdx}{\int Ydx}$ obtineat maximum minimumve valorem: existentibus $Z$ \& $Y$ functionibus quibuscunque ipsarum $x$, $y$, $p$, $q$, $r$, \&c. sive determinatis, sive indeterminatis.
\subsection*{Solutio.}
Casu quo sit $x=a$, sit $\int Zdx=A$, atque $\int Ydx=B$: eritque $A/B$ maximum vel minimum, siquidem relatio inter $x$ \& $y$ recte fuerit assignata. Erit igitur fractio $A/B$ æqualis eidem huic fractioni $\dfrac{\int Zdx}{\int Ydx}$, casu quo $x=a$, si alicubi una applicata $y$ augeatur particula $nv$. Tum vero fiet $\int Zdx$ æqualis ipsi $A$, una cum valore differentiali formulæ $\int Zdx$, qui sit $=dA$; similique modo $\int Ydx$ abibit in $B$ auctum valore differentiali formulæ $\int Ydx$, qui sit $=dB$; sicque ex adjecta particula $nv$ ad applicatam $y$, casu quo $x=a$, transibit fractio $\dfrac{\int Zdx}{\int Ydx}$ in hanc $\dfrac{A+dA}{B+dB}$; quæ æqualis esse debet fractioni $A/B$; \sourcepage{152}unde nascitur ista æquatio $BdA=AdB$; quæ præbebit æquationem inter $x$ \& $y$ quæsitam. Q. E. I.
\subsection*{Corollarium I.}
\originalsection{24}Ad hanc igitur æquationem inter $x$ \& $y$ inveniendam, effici debet ut valores differentiales ipsarum $\int Zdx$ \& $\int Ydx$ proportionales fiant ipsis harum formularum valoribus, quos obtinent posito $x=a$.
\subsection*{Corollarium II.}
\originalsection{25}Quanquam, in hac æquatione inventa $BdA=AdB$, duæ inesse videntur constantes incognitæ $A$ \& $B$, tamen ambas in unam compingere licet. Posito enim $A/B=C$, erit $dA=CdB$; inventaque æquatione, ex valore $a$ loco $x$ substituto determinabitur valor ipsius $C$.
\subsection*{Scholion.}
\originalsection{26}Si hujus \& præcedentis Problematis solutiones inter se conferantur, ingens in iis deprehendetur consensus. Nam si maximum minimumve esse debeat factum $\int Zdx\times\int Ydx$, orta est ista æquatio $0=AdB+BdA$; sin autem quotus $\dfrac{\int Zdx}{\int ydx}$ debeat esse vel maximus vel minimus, inventa est ista æquatio $0=AdB-BdA$; utroque autem casu litteræ $A$, $B$ \& $dA$, $dB$ eosdem retinent valores. Quare cum $A$ \& $B$ sint quantitates constantes, ambæ æquationes tantum ratione signi constantis differunt; posito enim $A/B=C$, priore casu habetur $dA=-CdB$, posteriore vero $dA=+CdB$. Ex quo pro utroque casu etiam eadem fere prodibit solutio; quia totum discrimen tantum in signo quantitatis constantis $C$ situm erit. Quod si ergo æquatio inter $x$ \& $y$ fuerit inventa, quæ
\sourcepage{153}contineat, pro $x=a$, factum $\int Zdx\times\int Ydx$ maximum vel minimum; eadem æquatio, levi adhibita mutatione, simul continebit quotum $\dfrac{\int Zdx}{\int Ydx}$ maximum vel minimum. Perspicuum autem est, sive $\dfrac{\int Zdx}{\int Ydx}$ debeat esse maximum vel minimum, sive $\dfrac{\int Ydx}{\int Zdx}$, utroque casu eandem plane esse prodituram æquationem. Hanc vero convenientiam ipsa rei natura postulat: nam si $\dfrac{\int Zdx}{\int Ydx}$ est maximum, tum eo ipso erit $\dfrac{\int Ydx}{\int Zdx}$ minimum \& vicissim; unde utrique quæstioni eandem solutionem satisfacere necesse est. Cæterum hunc quoque nexum observasse juvabit inter maximi minimive formulam $\dfrac{\int Zdx}{\int Ydx}$, quæ, posito $x=a$, abit in $A/B$, \& inter æquationem inventam $BdA-AdB=0$: hæc enim æquatio oritur ex differentiatione formulæ $A/B$, ponendo ejus differentiale $=0$; istiusmodi autem nexum perpetuo locum habere in sequente Propositione demonstrabimus.
\subsection*{Exemplum I.}
\originalsection{27}Invenire curvam, cujus area coordinatis orthogonalibus abscissa ad arcum curvæ maximam teneat rationem, si abscissæ datus valor $a$ tribuatur.

Posita curvæ quæsitæ abscissa $=x$, applicata $=y$; erit area $=\int ydx$, \& arcus $=\int dx\sqrt{1+pp}$; posito $dy=pdx$: maximum ergo esse debet $\dfrac{\int ydx}{\int dx\sqrt{1+pp}}$, casu quo ponitur $x=a$. Sit igitur, casu $x=a$, valor formulæ $\int ydx$, seu area $=A$, \& $\int dx\sqrt{1+pp}$ seu arcus abscissæ $a$ respondens $=B$. Deinde formulæ $\int ydx$ valor differentialis $dA$ erit $=nv.dx.1$, \sourcepage{154}\& formulæ $\int dx\sqrt{1+pp}$ seu
\[dB=nv.dx\left(-\frac{1}{dx}d.\frac{p}{\sqrt{1+pp}}\right)=-nv.d.\frac{p}{\sqrt{1+pp}}.\]
Quibus valoribus in æquatione $BdA=AdB$ substitutis, prodibit pro curva quæsita sequens æquatio: $Bdx=-Ad.\dfrac{p}{\sqrt{1+pp}}$. Ponatur $A/B=c$, ita ut, pro abscissa $x=a$, area curvæ fiat æqualis producto ex arcu in hanc constantem $c$. Erit ergo $dx=-cd.\dfrac{p}{\sqrt{1+pp}}$, \& integrando $x=b-\dfrac{cp}{\sqrt{1+pp}}$, seu $cp=(b-x)\sqrt{1+pp}$, hincque $p=\dfrac{b-x}{\sqrt{c^2-(b-x)^2}}=dy/dx$. Erit ergo
\[y=\int\frac{(b-x)dx}{\sqrt{c^2-(b-x)^2}}=f\pm\sqrt{c^2-(b-x)^2}\]
seu $(y-f)^2+(b-x)^2=cc$; unde constat curvam quæsitam esse Circulum radio $c$ descriptum, ad rectam quamcunque tanquam axem relatum. Hujus autem Circuli ea tantum portio quæsito satisfacit, quæ respondet abscissæ $=a$, a quo valore pendet $c$, ita ut sumpta abscissa $=a$, area æqualis fiat producto ex arcu in radium Circuli multiplicato. Quod si ergo vicissim radius $c$ detur, tanta in axe abscissa abscindi debet, ut arcus per radium multiplicatus præbeat aream. Infinitis igitur modis quæsito satisfieri potest; quæstio autem erit determinata, si duo præscribantur puncta, per quæ curva quæsita sit transeunda. Sumamus igitur radium $c$ tanquam cognitum, eoque describamus Circulum $BMD$ centro $C$.\marginnote{\footnotesize\hyperlink{fig:12}{Fig.~12.}} Porro sumatur linea quæcunque $APD$ pro axe, in eaque $A$ pro origine abscissarum. Hoc jam facto, quæstioni satisfiet si applicata $PM$ tantum spatium $ABMP$ abscindatur, ut id sit æquale producto ex arcu $BM$ in radium Circuli $BC$. Quia autem sector $BCM$ est $=\tfrac12BM.BC$, oportet aream $ABMP$ esse duplo majorem sectore $BCM$. Apparet autem, sumto pro lubitu, \sourcepage{155}cum axe, tum ejus initio, sæpe-numero conditionem præscriptam ne quidem impleri posse. Nam si axis $AD$ per centrum transeat, tum area $ABMP$ perpetuo minor erit quam duplum sectoris $BCM$; nisi, arcu $BM$ infinite parvo, prima applicata $BA$ simul per centrum transeat: sin autem axis $AD$ supra centrum transiret, tum nullo modo conditioni inventæ satisfieri potest. Quare necesse est, ut axis $AD$ infra centrum $C$ ducatur, qua de re multæ egregiæ observationes geometricæ fieri possent, si ratio instituti id permitteret. Cæterum si hæc Solutio cum Exemplo secundo præced. Prop. §. 20 comparetur; apparebit eandem prorsus æquationem esse inventam, sive $\int ydx\times\int dx\sqrt{1+pp}$ debeat esse minimum, sive $\dfrac{\int ydx}{\int dx\sqrt{1+pp}}$ maximum. Discrimen tamen in hoc consistit; quod radius Circuli $c=A/B$ altero casu affirmative, altero negative debeat accipi. Scilicet si $\int ydx\times\int dx\sqrt{1+pp}$ debeat esse minimum, arcus $BM$ convexitate sua spatium $ABMP$; altero autem casu, concavitate claudere debet.
\subsection*{Exemplum II.}
\originalsection{28}\marginnote{\footnotesize\hyperlink{fig:7}{Fig.~7.}}Intra datum angulum $ACM$, curvam $AM$ construere ita comparatam, ut area $ACM$ per arcum $AM$ divisa sit omnium maxima.

Ponatur angulus $ACM$, seu arcus circuli $BS$ radio $CB=1$ descriptus $=x$, qui in casu proposito fiat $=a$, quo $ACM/AM$ fieri debet maximum. Ponatur porro $CM=y$, sitque $dy=pdx$, erit $Mn=ydx$, \& area $ACM=\tfrac12\int yydx$: arcus autem $AM$ reperitur $=\int dx\sqrt{yy+pp}$: unde hæc fractio $\dfrac{\int yydx}{2\int dx\sqrt{yy+pp}}$, seu ejus duplum $\dfrac{\int yydx}{\int dx\sqrt{yy+pp}}$ debebit esse maximum. Sit, casu quo $x=a$ est, \sourcepage{156}$\int yydx=A$, \& $\int dx\sqrt{yy+pp}=B$; erit, si $x=a$, area $ACM=\tfrac12A$, \& arcus $AM=B$. Jam formulæ $\int yydx=A$ valor differentialis $dA$ est $=nv.dx.2y$, \& formulæ $\int dx\sqrt{yy+pp}$ valor differentialis $dB$ est
\[=nv.dx\left(\frac{y}{\sqrt{yy+pp}}-\frac{1}{dx}d.\frac{p}{\sqrt{yy+pp}}\right).\]
Quare cum generaliter invenerimus pro curva hanc æquationem $BdA=AdB$, erit divisione per $nv$ instituta,
\[2Bydx=\frac{Aydx}{\sqrt{yy+pp}}-Ad.\frac{p}{\sqrt{yy+pp}}.\]
Multiplicetur ea per $p$, ob $pdx=dy$, erit
\[2Bydy=A\left(\frac{ydy}{\sqrt{yy+pp}}-pd\frac{p}{\sqrt{yy+pp}}\right).\]
At est $d.\sqrt{yy+pp}=\dfrac{ydy}{\sqrt{yy+pp}}+\dfrac{pdp}{\sqrt{yy+pp}}$ \& $\dfrac{ydy}{\sqrt{yy+pp}}=d.\sqrt{yy+pp}-\dfrac{p}{\sqrt{yy+pp}}.dp$; unde fiet
\[2Bydy=A\left(d.\sqrt{yy+pp}-d.\frac{pp}{\sqrt{yy+pp}}\right);\]
ob $pd\dfrac{p}{\sqrt{yy+pp}}+dp\dfrac{p}{\sqrt{yy+pp}}=d.p.\dfrac{p}{\sqrt{yy+pp}}=d.\dfrac{pp}{\sqrt{yy+pp}}$. Quare integrando habebitur, si $A/B=c$ ponatur, ista æquatio
\[yy\pm bb=c\sqrt{yy+pp}-\frac{cpp}{\sqrt{yy+pp}}=\frac{cyy}{\sqrt{yy+pp}}\]
seu $p=\dfrac{y\sqrt{c^2y^2-(yy\pm bb)^2}}{yy\pm bb}=dy/dx$; hincque
\[dx=\frac{(yy\pm bb)dy}{y\sqrt{c^2y^2-(yy\pm bb)^2}}:\]
ex qua æquatione facile deduci potest, si sit $cc\mp4bb$ quantitas positiva, constructionem per quadraturam Circuli absolvi posse. At idem facilius patebit, si loco $dx$, vel $p$, introducamus perpendiculum $CP$, ex $C$ in tangentem $MP$ demissum. Quod si autem hoc perpendiculum $CP$ ponatur $=u$, erit $y:u=dx\sqrt{yy+pp}:ydx$, hincque $\dfrac{yy}{\sqrt{yy+pp}}=u$; quamobrem
\sourcepage{157}cum esset $yy\pm bb=\dfrac{cyy}{\sqrt{yy+pp}}$ erit $yy\pm bb=cu$, quam constat esse æquationem ad ipsum Circulum. Hoc ut ostendamus, sumatur Circulus quicunque, centro $O$, radio $OM=g$ descriptus, punctumque $C$ sumtum sit in $C$, ita ut sit $OC=h$.\marginnote{\footnotesize\hyperlink{fig:13}{Fig.~13.}} Jam ducta recta $CM=y$, \& $CP=u$, perpendiculari in tangentem $MP$, erit $CP$ parallela radio $OM$. Ex $M$ ducatur diametro $EF$ parallela $MR$, erit $MR=CO=h$; $CR=OM=g$, \& $PR=u-g$: quia igitur est $MR^2=MP^2+PR^2=CM^2-CP^2+PR^2$, erit $h^2=y^2-u^2+(u-g)^2=y^2-2gu+gg$; hincque $yy+gg-hh=2gu$; quæ comparata cum inventa $yy\pm bb=cu$, fiet $g=\tfrac12c$, \& $\pm bb=\tfrac14cc-hh$, seu $hh=\tfrac14cc\mp bb$. Erit itaque curva quæsita Circulus, radio $=\tfrac12c$ descriptus, puncto $C$ ubi libuerit accepto. In tali Circulo quæsito satisfaciet arcus $AM$, si fuerit $\dfrac{ACM}{AM}=\dfrac{A}{2B}=\tfrac12c=$ radio $OM$; hoc est si fuerit area $ACM=\mathrm{arc.}\ AM.AO=$ duplici sectori $AOM$. Hoc autem fieri nequit, nisi punctum $C$ extra Circulum accipiatur; quo casu hæc conditio infinitis modis adimpleri potest; atque adeo effici ut curva satisfaciens per data duo puncta transeat.
\subsection*{Exemplum III.}
\originalsection{29}\marginnote{\footnotesize\hyperlink{fig:14}{Fig.~14.}}Invenire curvam $DAD$ ad axem $AC$ relatam, in qua pro data abscissa $AC=a$, sit $\dfrac{\int xdx\sqrt{1+pp}}{\int dx\sqrt{1+pp}}$ minimum.

Si ponatur abscissa indefinita $AP=x$, applicata $PM=y$, \& $dy=pdx$; exprimit $\dfrac{\int xdx\sqrt{1+pp}}{\int dx\sqrt{1+pp}}$ distantiam centri gravitatis curvæ $MAM$, tanquam uniformiter gravis spectatæ a puncto infimo $A$; quæ ergo distantia, translato $P$ in $C$, debet esse minima. Ad hoc inveniendum, posito $x=a$, sit $\int xdx\sqrt{1+pp}=A$, \& $\int dx\sqrt{1+pp}=B$: formulæ autem $\int xdx\sqrt{1+pp}$ \sourcepage{158}reperitur valor differentialis $dA=-nv.d.\dfrac{xp}{\sqrt{1+pp}}$, \& formulæ $\int dx\sqrt{1+pp}$ valor differentialis $dB=-nv.d.\dfrac{p}{\sqrt{1+pp}}$; quibus in æquatione $BdA=AdB$ substitutis, prodibit
\[Bd.\frac{xp}{\sqrt{1+pp}}=Ad.\frac{p}{\sqrt{1+pp}},\]
\& posito $A/B=c$, erit $d.\dfrac{xp}{\sqrt{1+pp}}=cd.\dfrac{p}{\sqrt{1+pp}}$; unde integrando oritur $\dfrac{xp}{\sqrt{1+pp}}=\dfrac{cp}{\sqrt{1+pp}}-b$, seu $b\sqrt{1+pp}=(c-x)p$; hincque elicitur $p=\dfrac{b}{\sqrt{(c-x)^2-bb}}=dy/dx$. Erit ergo
\[y=\int\frac{bdx}{\sqrt{(c-x)^2-b^2}};\]
quæ æquatio indicat curvam quæsitam esse Catenariam, initio abscissarum pro $x$ in loco axis $AC$ quocunque accepto: quin etiam pro axe sumi potest recta quæcunque diametro Catenariæ $AC$ parallela, in eaque punctum quodcunque pro axis initio. Quomodocunque autem axis, ejusque initium constituatur, quæstioni satisfiet ea tantum curvæ portio, ubi sit $\int xdx\sqrt{1+pp}=c\int dx\sqrt{1+pp}$. Ponamus pro axe ipsam diametrum $AC$, \& verticem $A$ pro initio abscissarum accipi. Quia in $A$, ubi est $x=0$, sit $dy/dx=p=\infty$, necesse est ut sit $cc-bb=0$, ideoque $b=c$. Verum hoc casu fit $y=\int\dfrac{cdx}{\sqrt{xx-2cx}}$, quæ curva sursum directa fit imaginaria, donec fiat $x>2c$. Sit ergo $x=2c+t$, erit $t=$ abscissa $AP$, \& $y=PM=\int\dfrac{cdt}{\sqrt{2ct+tt}}$; curvaque $DAD$ erit catenaria ordinaria. Quo autem appareat quanta ejus portio quæstioni satisfaciat, notandum est, ob $dx=dt$, esse $p=\dfrac{c}{\sqrt{2ct+tt}}$, \& $\sqrt{1+pp}=\dfrac{c+t}{\sqrt{2ct+tt}}$; hincque
\[\int dx\sqrt{1+pp}=\int\frac{(c+t)dt}{\sqrt{2ct+tt}}=\sqrt{2ct+tt}.\]
At ipsa expressio $\dfrac{\int xdx\sqrt{1+pp}}{\int dx\sqrt{1+pp}}$ fit \sourcepage{159}$=2c+\dfrac{\int tdt\sqrt{1+pp}}{\sqrt{2ct+tt}}$, quæ ipsi $c$ æqualis fieri nullo modo potest. Ex quo concluditur nullam curvæ hujus portionem quæsito præ reliquis magis satisfacere. Quamobrem initium abscissarum non sumi potest in vertice $A$. Sumatur ergo in alio quocunque puncto, positaque $AP=t$, fieri debet $2bt+tt=(c-x)^2-bb$; unde fit, vel $b+t=x-c$, vel $b+t=c-x$. Prior æquatio $x=b+c+t$ locum habere nequit; quia, ob $dx=dt$, fieri non potest $\dfrac{\int xdt\sqrt{1+pp}}{\int dt\sqrt{1+pp}}$ seu $\left(b+c+\dfrac{\int tdt\sqrt{1+pp}}{\int dt\sqrt{1+pp}}\right)=c$. Ergo fiat $x=c-b-t$, quo casu abscissæ ab aliquo puncto axis $AC$ superiori deorsum descendent: fierique deberet $\dfrac{\int xdx\sqrt{1+pp}}{\int dx\sqrt{1+pp}}$ seu $c-b-\dfrac{\int tdt\sqrt{1+pp}}{\int dt\sqrt{1+pp}}=c$, quod pariter fieri nequit; ex quo concludendum est, nullam portionem magis quam aliam quamvis satisfacere. Hoc autem inde venire videtur, quod Catenaria duas habet partes conjugatas veluti Hyperbola conica, hincque semper fieri potest $\dfrac{\int xdx\sqrt{1+pp}}{\int dx\sqrt{1+pp}}=0$, qui est valor minimus. Hoc clarius confirmari potest ex valore invento $p=\dfrac{b}{\sqrt{(c-x)^2-b^2}}$; unde fit $\sqrt{1+pp}=\dfrac{c-x}{\sqrt{(c-x)^2-b^2}}=(c-x)r$, posito brevitatis gratia $r=\dfrac{1}{\sqrt{(c-x)^2-b^2}}$. Oporteret ergo in casu quæsito esse $\dfrac{\int(c-x)xrdx}{\int(c-x)rdx}=c$ seu $\int(c-x)^2rdx=0$, quod cum casu $x=0$ evanescere debeat, alio insuper casu evanescere deberet. At est
\[\begin{aligned}\int(c-x)^2rdx&=\int\frac{(c-x)^2dx}{\sqrt{(c-x)^2-b^2}}\\&=-\tfrac12(c-x)\sqrt{(c-x)^2-b^2}\\&\quad-\frac{bb}{2}l.\frac{c-x+\sqrt{(c-x)^2-b^2}}{c+\sqrt{c^2-b^2}}+\tfrac12c\sqrt{c^2-b^2},\end{aligned}\]
quæ expressio, cum semel fuit $=0$, post, ob $(c-x)^2$ perpetuo affirmativum, continuo crescet neque denuo \sourcepage{160}fieri potest $=0$. Quamobrem ambos terminos integrationis formulæ $\int\dfrac{(c-x)^2dx}{\sqrt{(c-x)^2-b^2}}$ inter se congruere oportet; quod evenit si fuerit $x=c$: quo casu curva satisfaciens abit in lineam rectam axi normalem, quæ utique centrum suum gravitatis a se minime habet remotum.
\subsection*{Exemplum IV.}
\originalsection{30}Invenire curvam, in qua, pro data abscissa $x=a$, sit hæc expressio $\dfrac{\int yxdx}{\int dx\sqrt{1+pp}}$ maximum vel minimum.

Posito $x=a$, fiet $\int yxdx=A$, \& $\int dx\sqrt{1+pp}=B$. Jam formulæ $\int yxdx$ valor differentialis est $dA=nv.dx.x=nv.xdx$, \& formulæ $\int dx\sqrt{1+pp}$ valor differentialis est $dB=-nv.d.\dfrac{p}{\sqrt{1+pp}}$. Quare cum sit $BdA=AdB$, habebitur ista æquatio $Bxdx=-d.\dfrac{p}{\sqrt{1+pp}}$, seu $xdx=-ccd.\dfrac{p}{\sqrt{1+pp}}$, posito $A=Bc^2$. Unde integrando obtinebitur $xx=C-\dfrac{2ccp}{\sqrt{1+pp}}$, hincque $p=\dfrac{bc-xx}{\sqrt{4c^4-(bc-xx)^2}}=dy/dx$; quæ præbet
\[y=\int\frac{(bc-xx)dx}{\sqrt{4c^4-(bc-xx)^2}},\]
quæ est æquatio generalis pro curva Elastica: cujus hæc proprietas, quod radius osculi ubique abscissæ $x$ sit reciproce proportionalis: id quod patet ex æquatione $xdx=-ccd.\dfrac{p}{\sqrt{1+pp}}$, quæ abit in $\dfrac{-dx}{d.\dfrac{p}{\sqrt{1+pp}}}=\dfrac{cc}{x}$, estque $\dfrac{-dx}{d.\dfrac{p}{\sqrt{1+pp}}}=-\dfrac{dx(1+pp)^{3:2}}{dp}$ radius osculi in curva. Hujus autem curvæ tanta portio ab initio computando satisfacit, in qua erit $\int yxdx$
\sourcepage{161}$=cc\int dx\sqrt{1+pp}=2c^4\int\dfrac{dx}{\sqrt{4c^4-(bc-xx)^2}}$; quæ determinatio eo revocatur, ut effici debeat
\[\int dx\sqrt{4c^4-(bc-xx)^2}=(aa-bc)\int\frac{dx(bc-xx)}{\sqrt{4c^4-(bc-xx)^2}},\]
si post integrationem utramque ponatur $x=a$. Hoc itaque modo constans illa $c$ per $a$ determinabitur.
\subsection*{Propositio IV. Problema.}
\originalsection{31}Invenire æquationem inter binas variabiles $x$ \& $y$ ita comparatam, ut, posita variabili $x=a$, maximum minimumve fiat expressio $W$, quæ sit functio quæcunque formularum integralium $\int Zdx$, $\int Ydx$, $\int Xdx$, \&c. in quibus denotent $Z$, $Y$, $X$ \&c. functiones quascunque ipsarum $x$, $y$, $p$, $q$, \&c. sive, determinatas, sive indeterminatas.
\subsection*{Solutio.}
Ponamus idoneam æquationem inter $x$ \& $y$ jam esse inventam, positoque $x=a$, fieri $\int Zdx=A$; $\int Ydx=B$; $\int Xdx=C$ \&c. hisque valoribus in expressione $W$ substitutis, habebitur revera maximum vel minimum. Quod si igitur altera variabilis $y$ in uno loco particula $nv$ augeri ponatur, atque nascentes hinc mutationes in singulis formulis $\int Zdx$, $\int Ydx$, $\int Xdx$ \&c. introducantur, idem pro $W$ valor prodire debet. At ab illa particula $nv$ formulæ $\int Zdx$, $\int Ydx$, \& $\int Xdx$ \&c. quæque suis valoribus differentialibus augebuntur. Si ergo ponatur formulæ $\int Zdx$ valor differentialis $=dA$, formulæ $\int Ydx=dB$, formulæ $\int Xdx=dC$, \&c. loco quantitatum $A$, $B$, $C$, \&c. orientur a particula $nv$ istæ auctæ $A+dA$, $B+dB$, $C+dC$ \&c. quæ in $W$ substitutæ eundem valorem producere debent, quem ipsæ $A$, $B$, $C$, \&c. Ponamus, $A+dA$, $B+dB$, $C+dC$ \&c. loco $\int Zdx$, $\int Ydx$, $\int Xdx$ \&c. substitutis, prodire $W+dW$; eritque $W+dW=W$, ideoque $dW=0$. Hic autem valor $dW$, ut ex differentiationis natura liquet, invenitur, \sourcepage{162}si quantitas $W$, postquam in illa, loco formularum integralium, litteræ $A$, $B$, $C$ \&c. sunt substitutæ, differentietur, his ipsis litteris $A$, $B$, $C$, \&c. tanquam variabilibus tractatis; in hocque differentiali, $dA$, $dB$, $dC$ \&c. valores differentiales formularum respondentium $\int Zdx$, $\int Ydx$, $\int Xdx$ \&c. designent. Hac igitur significatione sumptum differentiale quantitatis propositæ $W$, si id nihilo æquale ponatur, dabit æquationem inter $x$ \& $y$ quæsitam. Q. E. I.
\subsection*{Corollarium I.}
\originalsection{32}Si ergo proposita fuerit ejusmodi expressio $W$ functio formularum integralium $\int Zdx$, $\int Ydx$, $\int Xdx$ \&c. quæ, pro determinato ipsius $x$ valore $=a$, debeat esse maximum vel minimum: tum loco formularum $\int Zdx$, $\int Ydx$, $\int Xdx$ \&c. scribantur litteræ $A$, $B$, $C$, \&c. quo facto, expressio $W$ differentietur his litteris $A$, $B$, $C$, \&c. solis tanquam variabilibus tractatis, atque differentiale ponatur $=0$.
\subsection*{Corollarium II.}
\originalsection{33}In hoc differentiali, in quo inerunt litteræ $A$, $B$, $C$, \&c. cum suis differentialibus $dA$, $dB$, $dC$ \&c. litteræ $A$, $B$, $C$ \&c. denotabunt respective valores formularum $\int Zdx$, $\int Ydx$, $\int Xdx$ \&c. quos induunt posito $x=a$; at differentialia $dA$, $dB$, $dC$, \&c. exprimunt valores differentiales earundem formularum integralium abscissæ $x=a$ respondentes.
\subsection*{Corollarium III.}
\originalsection{34}Ex præcedentibus autem apparet, si $Z$, $Y$, $X$ \&c. fuerint functiones determinatæ quantitatum $x$, $y$, $p$, $q$, \&c. tum valores differentiales $dA$, $dB$, $dC$, \&c. non a valore $a$ pendere: contra vero si $Z$, $Y$, $X$ \&c. fuerint functiones indefinitæ, tum valores differentiales $dA$, $dB$, $dC$ \&c. simul a valore $a$ pendere debere.
\subsection*{Corollarium IV.}
\originalsection{35}\sourcepage{163}Cum igitur hoc modo $W$ fiat functio litterarum $A$, $B$, $C$, \&c. ejus differentiale hujusmodi habebit formam $FdA+GdB+HdC+\&c.$ hincque æquatio quæsita erit $0=FdA+GdB+HdC+\&c.$ ubi $F$, $G$, $H$ \&c. erunt quantitates constantes, per $A$, $B$, $C$ \&c. determinatæ.
\subsection*{Corollarium V.}
\originalsection{36}Æquatio ergo Problemati satisfaciens, constabit ex valoribus differentialibus singularum formularum integralium in maximi minimive expressione $W$ contentarum, singulis per constantes quantitates determinatas multiplicatis: horum scilicet productorum aggregatum nihilo æquale positum dabit æquationem desideratam.
\subsection*{Scholion I.}
\originalsection{37}Potuissemus hanc Problema propositum resolvendi methodum, ex solutionibus binorum Problematum præcedentium, per inductionem, jam concludere: quippe ex quibus jam patebat, si fuerit maximi minimive formula $W$, vel productum ex duabus formulis integralibus, vel quotus ex divisione unius per alteram ortus, tum differentiale expressionis $W$ modo exposito sumtum præbere æquationem Problemati convenientem. Præstitit autem hoc Problema, ob summam ejus extensionem, singulari solutione munire. In hoc enim Problemate continentur omnes omnino Quæstiones, quæ in hoc genere, quo expressio quæpiam maxima minimave desideratur, unquam proponi atque excogitari possunt: ideoque per istam Propositionem penitus exhausta est methodus maximorum ac minimorum absoluta, quam primo pertractandam suscepimus. Præterea hic notandum est, si expressio $W$ non tantum formulas integrales, uti posuimus, \sourcepage{164}complectatur; verum etiam functiones determinatas ipsarum $x$, $y$, $p$, $q$, \&c. tum solutionem nihilo difficiliorem reddi. Nam pari modo, loco harum functionum determinatarum, quantitates constantes poni debent, in quas scilicet abeunt posito $x=a$; at postmodum, in differentiatione ipsius $W$, has quantitates etiam tanquam constantes tractari oportet; eo quod functiones determinatæ nullos valores differentiales recipiunt. Quo autem clarius appareat, quomodo istiusmodi expressiones tractari conveniat; in sequentibus Exemplis nonnulla occurrent, quæ hoc argumentum penitus illustrabunt.
\subsection*{Exemplum I.}
\originalsection{38}Invenire curvam coordinatis orthogonalibus contentam, in qua sit maximum vel minimum ista expressio $(1+pp)^{1:2}\int ydx+y\int dx\sqrt{1+pp}$; si ponatur abscissa $x=a$.

Ponamus æquationem inter $x$ \& $y$ quæsito satisfacientem jam esse inventam, atque posito $x=a$ fieri $y=f$ \& $\sqrt{1+pp}=g$: itemque $\int ydx=A$, \& $\int dx\sqrt{1+pp}=B$; erit $dA=nv.dx$, \& $dB=-nv.d\dfrac{p}{\sqrt{1+pp}}$. Expressio igitur, quæ maxima erit vel minima, hoc casu est $gA+fB$, cujus differentiale est $gdA+fdB$; quod positum $=0$, dabit æquationem desideratam pro curva. Hic scilicet intelligitur litteras $y$ \& $f$, quæ ex functionibus determinatis sunt ortæ, in differentiatione tanquam quantitates constantes esse tractatas. Substitutis jam pro $dA$ \& $dB$ valoribus debitis, divisioneque per $nv$ facta, orietur ista æquatio pro curva quæsita $gdx=fd.\dfrac{p}{\sqrt{1+pp}}$. Ponatur $f/g=c$, ita ut sit $\dfrac{y}{\sqrt{1+pp}}=c$, casu quo est $x=a$; erit integrando $x+b=\dfrac{cp}{\sqrt{1+pp}}$, atque $p=$
\sourcepage{165}$\dfrac{x+b}{\sqrt{cc-(x+b)^2}}=dy/dx$; ex qua fit $y=h\pm\sqrt{c^2-(x+b)^2}$. Curva igitur satisfaciens est Circulus, radio $c$ descriptus, abscissis super recta quacunque assumptis, pariterque abscissarum initio ubicunque statuto. Quantitas autem $c$, quæ radium Circuli constituit, ex definita abscissa $x=a$ determinatur; quia esse debet $\dfrac{y}{\sqrt{1+pp}}=c$, casu quo $x=a$. Fit autem hoc casu $y=h\pm\sqrt{c^2-(a+b)^2}$, \& $\sqrt{1+pp}=\dfrac{c}{\sqrt{cc-(a+b)^2}}$: unde oritur $cc=h\sqrt{cc-(a+b)^2}\pm(cc-(a+b)^2)$, per quam, vel $c$ per $a$, vel vicissim $a$ per $c$ determinari potest. Ponamus esse $h=0$, $b=-c$, ita ut axis sit Circuli diameter, initiumque abscissarum in vertice constituatur; erit $y=\sqrt{2cx-xx}$, atque fiet $(a-c)^2=0$, seu $c=a$. Ex quo intelligitur, hoc casu quadrantem Circuli quæsito satisfacere. Sin autem initium abscissarum in loco diametri quocunque capiatur, fiet tantum $h=0$, \& si applicatæ positivæ sumantur fiet $(a+b)^2=0$, seu $b=-a$. Diameter Circuli ergo manet indeterminatus: portioque Circuli hoc modo sumti quæstioni satisfaciet, quæ abscissæ a sua origine ad centrum Circuli usque productæ respondet.
\subsection*{Exemplum II.}
\originalsection{39}Invenire æquationem inter $x$ \& $y$, ut pro valore definito $x=a$, hæc expressio $y^{\int dx\sqrt{1+pp}}\int ydx$ fiat maximum vel minimum.

Posito $x=a$, fiat $y=f$, $\int dx\sqrt{1+pp}=A$, \& $\int ydx=B$, erit $dA=-nv.d.\dfrac{p}{\sqrt{1+pp}}$ \& $dB=nv.dx$. Maximum ergo minimumve esse oportet hanc quantitatem $f^AB$, cujus differentiale est $f^ABdAlf+f^AdB$; quod positum $=0$ dabit $BdAlf=-dB$. Pro æquatione quæsita igitur habetur \sourcepage{166}$Blfd.\dfrac{p}{\sqrt{1+pp}}=dx$, \& integrando $x+b=\dfrac{Bplf}{\sqrt{1+pp}}=\dfrac{cp}{\sqrt{1+pp}}$, posito $Blf=c$. Habetur ergo $p=\dfrac{b+x}{\sqrt{cc-(b+x)^2}}$ \& $y=h\pm\sqrt{c^2-(b+x)^2}$. Erit igitur $f=h\pm\sqrt{c^2-(b+a)^2}$, posito $x=a$, atque $B=\int ydx=ha\pm\int dx\sqrt{c^2-(b+x)^2}$, posito post integrationem $x=a$. Facto igitur $Blf=c$, innotescet valor $a$, cui si $x$ æqualis capiatur in Circulo radii $c$, portio abscindetur Problemati satisfaciens. Cæterum ex his \& Coroll. 5 colligere licet, quoties formula maximi minimive fuerit functio quæcunque binarum harum formularum $\int ydx$ \& $\int dx\sqrt{1+pp}$, curvam satisfacientem perpetuo esse Circulum: tantum ex solutione quantitas portionis satisfacientis debet diligenter investigari ac determinari.
\subsection*{Exemplum III.}
\originalsection{40}Invenire æquationem inter $x$ \& $y$, ut, posito $x=a$, maximum minimumve fiat ista expressio
\[e^{-n\int dx\sqrt{1+pp}}\int e^{n\int dx\sqrt{1+pp}}dx.\]
Ponamus, casu proposito quo $x=a$, fieri $n\int dx\sqrt{1+pp}=A$, atque $\int e^{n\int dx\sqrt{1+pp}}dx=B$; ita ut maximum minimumve sit hæc quantitas $e^{-A}B$, cujus differentiale est $e^{-A}dB-e^{-A}BdA$; quod positum $=0$ dabit æquationem hanc $dB=BdA$. At est $dA$ valor differentialis formulæ $n\int dx\sqrt{1+pp}$, unde erit $dA=-nv.d.\dfrac{np}{\sqrt{1+pp}}$: atque $dB$ est valor differentialis formulæ $\int e^{n\int dx\sqrt{1+pp}}dx$, quæ continetur in Casu secundo §. 7, ubi est $Z=e^{n\int dx\sqrt{1+pp}}$ \& $\Pi=\int dx\sqrt{1+pp}$, ita ut sit $Z=e^{n\Pi}$, \& $dZ=e^{n\Pi}nd\Pi$, unde erit $L=e^{n\Pi}n$, \& reliquæ litteræ $M$, $N$, \sourcepage{167}$P$, \&c. fient $=0$. Porro, ob $\Pi=\int dx\sqrt{1+pp}$, erit $[Z]=\sqrt{1+pp}$, \& $d[Z]=\dfrac{pdp}{\sqrt{1+pp}}$, ex quo erit $[M]=0$, $[N]=0$, \& $[P]=\dfrac{p}{\sqrt{1+pp}}$. Jam est $\int Ldx=n\int e^{n\int dx\sqrt{1+pp}}dx$, cujus valor, posito $x=a$, erit $=nB$; hincque $V=n\left(B-\int e^{n\int dx\sqrt{1+pp}}dx\right)$. Per Regulam ergo datam fiet
\[\begin{aligned}dB&=nv.dx\left(-\frac{d.[P]V}{dx}\right)\\&=-nv.d.\frac{np(B-\int e^{n\int dx\sqrt{1+pp}}dx)}{\sqrt{1+pp}}\\&=-nv.d.\frac{nBp}{\sqrt{1+pp}},\end{aligned}\]
ob $dB=BdA$. Integrando itaque erit
\[\frac{np(B-\int e^{n\int dx\sqrt{1+pp}}dx)}{\sqrt{1+pp}}=\frac{nBp}{\sqrt{1+pp}}-nb;\]
hincque $\dfrac{b\sqrt{1+pp}}{p}=\int e^{n\int dx\sqrt{1+pp}}dx$. Ex qua æquatione, quia valor determinatus $a$ excessit, perspicuum est æquationem inventam pro quovis ipsius $x$ valore æque valere. Ut autem hanc æquationem evolvamus, erit, differentialibus sumtis,
\[\frac{-bdp}{p^2\sqrt{1+pp}}=e^{n\int dx\sqrt{1+pp}}dx:\]
quæ, per $\sqrt{1+pp}$ multiplicata atque integrata, dat $nb/p+c=e^{n\int dx\sqrt{1+pp}}$, qui exponentialis quantitatis valor in illa æquatione substitutus, dabit
\[\frac{nbdx}{p}+cdx=-\frac{bdp}{p^2\sqrt{1+pp}},\qquad\text{seu}\qquad dx=\frac{-bdp}{p(nb+cp)\sqrt{1+pp}}.\]
Commodior autem æquatio oritur, si ponatur $\int dx\sqrt{1+pp}=s$, eritque $s$ arcus curvæ, si fuerint $x$ \& $y$ coordinatæ normales. Quare habebitur ista æquatio $nb+cp=e^{ns}p$, quæ per $dx$ multiplicata, ob $dy=pdx$, abit in hanc $nbdx+cdy=e^{ns}dy$.

\sourcepage{168}Cum autem, posito, $x=0$, arcus $s$ evanescere debeat, necesse est ut sit hoc casu $nb/p+c=0$; hinc itaque vel dato curvæ initio constans $c$ determinabitur, vel vicissim ex $c$ positio primæ tangentis innotescet. Cæterum si hanc quæstionem attentius contemplemur, deprehendemus eam jam contineri in Exemplo quodam Capitis præced. §. 45. Cum enim nostra expressio, quæ maximum minimumve esse debeat, sit $e^{-n\int dx\sqrt{1+pp}}\int e^{n\int dx\sqrt{1+pp}}dx$; ponatur ea $=W$, erit $e^{n\int dx\sqrt{1+pp}}W=\int e^{n\int dx\sqrt{1+pp}}dx$, atque differentiando fiet $dW+nWdx\sqrt{1+pp}=dx$. Maximi igitur minimive expressio $W$ datur per æquationem differentialem, quæ in Casu quarto §. 7 continetur: atque methodo convenienti tractata ad eandem perducit æquationem, quam hic invenimus. Quæstionem autem illam in se complectentem supra in Cap. præc. §. 45 tractavimus, in quo hunc ipsum casum adjunctum spectare licet. Comparatione autem instituta, summus perspicietur consensus solutionum variarum ejusdem Problematis, quæ quidem tentari queant.
\subsection*{Exemplum IV.}
\originalsection{41}Invenire curvam in qua, pro data abscissa $=a$, fiat ista expressio
\[\frac{\int dx\sin.\mathrm{A}.y.\sqrt{1+pp}}{\int dx\cos.\mathrm{A}y.\sqrt{1+pp}}\]
maximum vel minimum.

Posito $x=a$, fiat $\int dx(1+pp)^{1:2}\sin.\mathrm{A}y=A$, \& $\int dx(1+pp)^{\frac12}\cos.\mathrm{A}y=B$; erit, per valores differentiales,
\[\begin{aligned}dA&=nv.dx\left((1+pp)^{1:2}\cos.\mathrm{A}y-\frac{1}{dx}d.\frac{p\sin.\mathrm{A}.y}{\sqrt{1+pp}}\right),\\dB&=nv.dx\left(-(1+pp)^{1:2}\sin.\mathrm{A}y-\frac{1}{dx}d.\frac{p\cos.\mathrm{A}y}{\sqrt{1+pp}}\right).\end{aligned}\]
Cum igitur $A/B$ debeat esse maximum vel minimum, erit $BdA$
\sourcepage{169}$=AdB$; posito ergo $A/B=m$, fiet
\[\begin{aligned}(1+pp)^{1:2}dx\cos.\mathrm{A}y-d.\frac{p\sin.\mathrm{A}y}{\sqrt{1+pp}}&=-m(1+pp)^{\frac12}dx\sin.\mathrm{A}y\\&\quad-md.\frac{p\cos.\mathrm{A}y}{\sqrt{1+pp}}.\end{aligned}\]
Multiplicetur per $p$, erit [ob
\[\begin{aligned}d.(1+pp)^{1:2}\sin.\mathrm{A}y&=dy(1+pp)^{1:2}\cos.\mathrm{A}y+\frac{pdp\sin.\mathrm{A}y}{\sqrt{1+pp}},\\d.(1+pp)^{1:2}\cos.\mathrm{A}y&=-dy(1+pp)^{1:2}\sin.\mathrm{A}.y+\frac{pdp\cos.\mathrm{A}y}{\sqrt{1+pp}}\end{aligned}\]
]
\[\begin{aligned}d.(1+pp)^{1:2}\sin.\mathrm{A}y-d.\frac{pp\sin.\mathrm{A}y}{\sqrt{1+pp}}&=md.(1+pp)^{1:2}\cos.\mathrm{A}y\\&\quad-md.\frac{pp\cos.\mathrm{A}y}{\sqrt{1+pp}};\end{aligned}\]
quæ integrata \& reducta præbet $\dfrac{\sin.\mathrm{A}y}{\sqrt{1+pp}}=\dfrac{m\cos.\mathrm{A}y}{\sqrt{1+pp}}+b$, sive $b\sqrt{1+pp}=\sin.\mathrm{A}y-m\cos.\mathrm{A}y$; ubi notandum est fieri debere, si $x=a$ ponatur,
\[m=\frac{\int dx(1+pp)^{1:2}\sin.\mathrm{A}y}{\int dx(1+pp)^{1:2}\cos.\mathrm{A}y}.\]
Sit $m=\dfrac{\sin.\mathrm{A}n}{\cos.\mathrm{A}n}=\mathrm{tang.}\ \mathrm{A}n$; fiet $b\sqrt{1+pp}=\dfrac{\sin\mathrm{A}(y-n)}{\cos.\mathrm{A}n}$, atque $y=n+\mathrm{A\ sin.}b(1+pp)^{1:2}\cos.\mathrm{A}n$. Quia vero est $dy=pdx$, erit $dx=dy/p$. At est $dy=\dfrac{cpdp}{\sqrt{(1+pp)(1-cc-ccpp)}}$, posito $b\cos.\mathrm{A}n=c$. Ex quibus conficitur
\[x=\int\frac{cdp}{\sqrt{(1+pp)(1-cc-ccpp)}},\qquad\text{atque}\qquad y=\int\frac{cpdp}{\sqrt{(1+pp)(1-cc-ccpp)}};\]
longitudo autem curvæ erit $=\int\dfrac{cdp}{\sqrt{1-cc-ccpp}}=\mathrm{A\ sin.}\dfrac{cp}{\sqrt{1-cc}}$. Quare si arcus curvæ dicatur $s$; habebitur ista concinna æquatio $dx\sin\mathrm{A}s=\dfrac{cdy}{\sqrt{1-cc}}$: Constructio vero ex anterioribus formulis sponte consequitur.
\subsection*{Scholion II.}
\originalsection{42}\sourcepage{170}His igitur Capitibus penitus absolvimus eam Methodi maximorum ac minimorum ad lineas curvas inveniendas accommodatæ partem, quam absolutam vocavimus: in qua semper linea curva requiri solet, quæ habeat, pro dato quodam abscissæ seu alterius variabilis $x$ valore, expressionem quamcunque indeterminatam, maximum minimumve. Nam ista expressio, quæ maximum minimumve esse debet, vel erit una quædam formula integralis formæ $\int Zdx$, ita ut $Z$ sit functio quæcunque ipsarum $x$, $y$, $p$, $q$, \&c. sive definita sive indefinita; pro quibus casibus Methodum tradidimus in Capitibus præcedentibus: Vel maximi minimive expressio illa continebit in se plures ejusmodi formulas integrales, ita ut sit duarum pluriumve formularum integralium functio quæcunque; pro hocque casu Methodus idonea in isto Capite est exposita, atque Exemplis illustrata. Universa autem Methodus, quam hic dedimus, nititur inventione valorum differentialium, qui singulis formulis integralibus quæ vel ipsæ maximum minimumve esse debeant, vel in maximi minimive expressione contineantur, atque ideo tota solvendi Methodus reducitur ad Casus illos, quos §. 7 hujus Capitis conjunctim repræsentavimus. Qui igitur illos casus in memoria tenet, vel in promtu habet, is ad omnia hujus generis Problemata expedite resolvenda erit paratus. Neque vero solum Casus ibi enumerati Methodum maximorum ac minimorum absolutam constituunt; verum etiam Methodum alteram relativam, quam in sequentibus aggrediemur, absolvent; ex quo illorum Casuum summus usus in utraque Methodo abunde perspicietur. Hanc autem tractationem duobus Capitibus absolvemus, in quorum priori omnibus curvis, ex quibus quæsita debet erui, unam quandam proprietatem communem, in posteriori vero plures tribuemus.
